% Contents/lutas_e_combates.tex
\midsectiontitle{Combate em OverGrown}

\titlesection{A Rolagem de Iniciativa}

No \textit{Overgrown}, o combate não é um mero caos, mas uma sequência de movimentos calculados onde a velocidade e a percepção ditam a sobrevivência. Quando as lâminas se cruzam, a ordem das ações é determinada pelo teste de \link{att:grace}{Grace} somado ao Modificador de Iniciativa. Este valor representa a prontidão do aventureiro em resposta à ameaça iminente.

\begin{rpgboxtips}{Rolagem de Iniciativa}
A iniciativa é rolada no início de cada combate com base em \link{att:grace}{Grace}. Aquele que obtiver o maior resultado age primeiro, seguido pelos demais em ordem decrescente. Em caso de empate, o personagem com maior \link{att:grace}{Grace} age primeiro; persistindo o empate, o Mestre decide ou permite nova rolagem.
\end{rpgboxtips}

Aqueles que se destacam agem primeiro, estabelecendo o ritmo do confronto, seguidos pelos demais em ordem decrescente de seus resultados.

\begin{rpgboxtips}{A Estratégia do Atraso}
Um combatente astuto sabe que nem sempre golpear primeiro é a melhor escolha. Você pode optar por \textbf{atrasar sua ação}, aguardando o momento exato em que um aliado abre uma brecha ou um inimigo se expõe. Ao fazer isso, sua posição na ordem de combate é permanentemente redefinida para este novo marco até o desfecho da cena.
\end{rpgboxtips}

\titlesection{O Fluxo do Tempo e Ações}

O tempo em combate é uma percepção fragmentada. Cada turno condensa entre 0,5 e 6 segundos, onde decisões devem ser tomadas com precisão antes de serem enunciadas. Sob o olhar do mestre, o impossível pode ser tentado, mas apenas o que for coerente com o fluxo temporal será executado.

\subsection*{A Hierarquia das Ações}

A economia de combate em \textit{Overgrown} é baseada na versatilidade. Cada combatente possui uma quantidade limitada de energia e pontos de combate que se traduz em quatro tipos de ações: \textbf{Padrão (ou Completa)}, \textbf{Movimento}, \textbf{Livre} e \textbf{Reação}.

A maestria reside na conversão de esforços. O sistema permite que o jogador sacrifique uma ação de maior complexidade para realizar múltiplas tarefas mais simples.

\begin{itemize}
    \item \textbf{Ação Padrão:} O núcleo do seu turno. Usada para atacar ou usar magias. Pode ser convertida em uma ação de Movimento, Livre ou em uma Reação adicional.
    \item \textbf{Ação de Movimento:} Utilizada para reposicionamento. Pode ser convertida em uma ação Livre.
    \item \textbf{Ação Livre:} Ações rápidas ou frases curtas. Não pode ser convertida.
    \item \textbf{Reação:} Uma resposta instintiva fora do seu turno (1 por rodada). Não pode ser convertida ou trocada. Esquivar de um ataque, bloquear, contra-atacar e interromper conjurações são exemplos de Reações.
\end{itemize}

\begin{rpgboxtips}{Conversão em Reação}
Uma Ação Padrão só pode ser convertida em Reação durante o próprio turno. A Reação adicional fica disponível até o início do próximo turno do personagem e desaparece caso não seja usada. Ela se soma à Reação normal da rodada, permitindo no máximo duas Reações nesse intervalo, salvo quando uma habilidade declarar outro limite. A conversão não pode ser feita retroativamente depois que já ocorreu.
\end{rpgboxtips}

\titlesection{Nível de Ameaça e Precisão}

\subsection*{Teste de Ataque}

Para resolver um ataque, role a reserva de d20 determinada pelo atributo indicado pela arma, magia ou habilidade e conserve o maior resultado. Em seguida, some a Perícia de Combate apropriada e quaisquer bônus ou penalidades que declarem afetar o teste de ataque:

\begin{itemize}
    \item \textbf{Luta:} ataques corpo a corpo, incluindo armas brancas e ataques desarmados.
    \item \textbf{Pontaria:} ataques com armas de longo alcance, mesmo quando a arma utilizar MIG em vez de GRA para formar sua reserva.
    \item \textbf{Arcanismo:} ataques de magias e ataques básicos de armas mágicas que utilizem WIS.
\end{itemize}

\[\text{Resultado de Ataque}=\text{maior d20 da reserva}+\text{Perícia de Combate}+\text{bônus e penalidades de acerto}\]

O ataque acerta quando seu resultado iguala ou supera a defesa aplicável do alvo. Um 20 natural continua sendo um acerto crítico garantido e um 1 natural continua sendo uma falha crítica, independentemente do total após os bônus.

O \textbf{MOD de Atributo não participa do teste de ataque}. Expressões como \textbf{MOD de MIG}, \textbf{MOD de GRA}, \textbf{MOD de WIS} ou \textbf{MOD de SEN} presentes na ficha da arma ou da magia são somadas somente ao dano, depois de confirmar o acerto. Por exemplo, uma arma corpo a corpo com dano \textbf{2D6 + MOD de MIG} testa o acerto usando a reserva indicada + Luta; ao acertar, causa 2D6 + MOD de MIG. Uma arma física à distância testa com Pontaria e acrescenta MOD de SEN ao dano.

\subsection*{Regra unificada de ataques à distância}
Arcos, bestas, pistolas, rifles, espingardas e demais armas físicas de longo alcance usam \textbf{Pontaria} no teste de ataque e \textbf{MOD de SEN} no dano. Exemplo: um Rifle de Precisão causa \textbf{2D10 + MOD de SEN}. Uma arma híbrida usa SEN apenas no modo à distância; uma Lança arremessada usa SEN, mas empunhada em corpo a corpo mantém MIG. Magias projetadas e focos arcanos continuam usando a regra da magia e seu atributo mágico.

\originlink{stats:crit}
No mundo de \textit{Overgrown}, o \textbf{Nível de Ameaça} descreve o potencial do aventureiro de encontrar brechas na defesa inimiga e atingir pontos vitais com precisão absoluta.

Por padrão, um Crítico ocorre ao obter um \textbf{20 natural} nos dados. Entretanto, através de armas específicas e do avanço nos conhecimentos da \textbf{Árvore de Talentos}, este limiar de precisão pode ser reduzido. Um aventureiro de elite pode expandir sua margem de ameaça para resultados como 19, 18 ou até menores, tornando cada golpe uma ameaça potencialmente letal.

\begin{rpgboxtips}{Nós de Agilidade}
Não ignore a Árvore de Talentos durante seus treinos. Existem \textbf{Nós de Estatísticas} e \textbf{Nós Supremos} focados especificamente em quebrar as regras de tempo, permitindo que ações de movimento concedam bônus extras ou que reações sejam recuperadas mais rapidamente.
\end{rpgboxtips}

Um aventureiro por padrão tem 7 unidades de movimento, mas para cada 10 pontos de Grace, ele ganha 1 unidade adicional. Assim, um personagem com 25 de Graça teria 9 unidades de movimento, permitindo-lhe se deslocar mais rapidamente pelo campo de batalha e alcançar posições estratégicas com maior facilidade.

\titlesection{Cobertura}

Obstáculos físicos entre o atirador e o alvo concedem proteção. Em \textit{Overgrown}, a cobertura é classificada em dois graus:

\begin{description}
    \item[\ringm{Cobertura Parcial}:] O alvo está protegido por um obstáculo que cobre \textbf{metade ou menos} de seu corpo (parapeito, pilar, canto de parede). O atirador sofre \textbf{$-$5 de acerto} nos ataques à distância contra esse alvo. O alvo ainda pode ser atingido normalmente caso o ataque supere essa penalidade.
    \item[\ringm{Cobertura Total}:] O alvo está completamente atrás de um obstáculo sólido, sem linha de visão direta. Ataques à distância comuns são \textbf{impossíveis} contra alvos em cobertura total. Apenas habilidades específicas (como \textit{Tiro Ricochete} ou \textit{Tiro em Curva}) ou magias de área permitem contornar este impedimento.
\end{description}

\begin{rpgboxtips}{Cobertura e Corpo a Corpo}
Cobertura não oferece proteção contra ataques corpo a corpo. Um inimigo que atravessar o obstáculo ou alcançar o alvo fisicamente ignora qualquer bônus de cobertura.
\end{rpgboxtips}

\titlesection{Canalização Mágica com Armas}

Certas armas possuem afinidade com energias arcanas e podem ser usadas como catalisadores para magias, transformando projéteis comuns em veículos de destruição mágica.

\begin{description}
    \item[\ringm{Arcos}:] Arcos são excelentes condutores de energia primordial. Ao canalizar uma magia através de um arco, o conjurador dispara a magia como se fosse uma flecha. O ataque utiliza o \link{stats:crit}{Nível de Ameaça} do arco e seus bônus de distância. A magia herda o tipo de dano definido pelo conjurador. Por continuar sendo uma magia, o teste soma \textbf{Arcanismo} e mantém \textbf{Wisdom}; não se converte em um ataque físico de Sense.
    \item[\ringm{Armas de Fogo}:] Por serem forjadas com gravuras de combustão rúnica focadas em dano físico bruto, armas de fogo \textbf{não podem ser usadas como canalizadores mágicos}. Tentar fazê-lo dissipa a magia sem efeito e gasta o IEP normalmente.
\end{description}

\begin{rpgboxtips}{Flechas Especiais e Magia}
Flechas de Aetherium e Rúnicas amplificam ainda mais magias canalizadas. O dano extra da flecha é adicionado ao dano da magia disparada, tornando a combinação de arqueria e arcanismo uma especialização poderosa.
\end{rpgboxtips}
